\section{The source raw-theta supremum estimate with constant \(9000\)}
\label{sec:sr-source-raw-sup-9000}

This section translates \srcid{RawThetaSup.homN_raw_le} and
\srcid{homN_smul_raw_le}.  The absolute constant is chosen before
the host, cover, weight, lattice, or scalar.  The proof uses the
orthogonal half-maximum and raw-value estimates of
Sections~\ref{sec:sr-source-raw-theta-norm} and
\ref{sec:sr-source-half-maximum}.  Their actual analytic dictionaries
remain the same dictionaries throughout the following argument.

\subsection{The exact family and its hypotheses}
Use the source packet \(D:\mathrm{D7}\), \(T:\mathrm{CoreJ}(D)\),
\(j\in\mathbb N\), and \(Y:\mathrm{Deep}(T,j)\).
Write
\[
 E=D.\mathrm{host}.E,\quad f=D.R.f=\tfrac12[E:\mathbb Q],\quad
 r=D.R.r,\quad D_w=T.Q.Dw(j),\quad M=Y.d.l_2.M,\quad
 c=Y.d.c_2 .
\]
The actual packet satisfies \(f,r\ge1\).
For a nonnegative exponent array \(a\) and even \(n\), put
\[
 d=\mathrm{DF}(D_w,a),\qquad
 \kappa=\mathrm{sizeK}(D_w,a,n)
             =d+\frac n2fr,\qquad \xi=\xi_R(D)>0.
 \tag{RN1}\label{eq:sr-source-raw-size}
\]
The equality for \(\kappa\) will be proved below, rather than
assumed as an independent normalization.

We retain every clause of the source
\(\mathrm{CovGmin.Hyp}(D,T,j,Y,e,o,L_b,a,n)\):
\begin{enumerate}
\item \(\mathrm{LawsJ}(T)\), \(D_w.p,D_w.q\ge1\),
\(2<[E:\mathbb Q]/2\), and \(\mathrm{InR}(D,j)\);
\item the actual \(\mathrm{Good}(Y,e,o)\) datum;
\item \(L_b\ge1\),
\(\log M\le30L_b\), \(\mathrm{fieldSz}(E)\le3L_b\),
\(\log(fr)\le L_b\), and \(o\le391(frL_b)^2\);
\item \(\mathrm{Dom}(a,n)\), \(\mathrm{Even}(n)\),
\(8fr+4\le n\), and
\[
  n/2<\mathrm{laOf}(D_w,a,n;\tau,k)
       \quad(\tau\in\Psi,\ k<D_w.p);
\]
\item \(\log\kappa\le20L_b\).
\end{enumerate}
Here \(\mathrm{Dom}\) is the source's dominant exponent condition,
not a replacement by positivity of total weight.
The separate evenness clause and the strict row inequalities are
kept as they occur in the source.

Fix an actual \(\mathrm{nab}\) satisfying \(\mathrm{IsNab}\) for
\(Y.d.Xo_2,c\).  Let \(x_0\) be a tuple in
\(\mathrm{LevQc}(D_w,l_2)\), and let \(\Lambda\) be the actual
positive coefficient lattice satisfying
\[
 \mathrm{Unimod}(\Lambda),\qquad
 \mathrm{TrivAt}(\Lambda,M),\qquad
 \mathrm{LevOK}(\Lambda,N,M).
\]
Let \(P\in\mathrm{harmL}(c,\Lambda,a)\), with
\(\mathrm{polyNorm}_c(P)\le\rho\).
Let \(y\in Y.V(\mathrm{wtZ}(a,n))\), and assume the exact
\(\mathrm{thLaValAt}\) formula for \(y\), for every source test ring
and every frame over that test ring, using these same
\[
 \mathrm{nab},\quad
 w=\phiHat_c(\mathrm{indPhi}_c(x_0)),\quad \Lambda,N,P .
 \tag{RN2}\label{eq:sr-source-raw-all-tests}
\]
In particular \(y\) is in the actual cusp space.  The estimate
does not replace (RN2) by equality at one complex frame.

Set
\[
 J_\Lambda=
 [\Lambda:\Lambda\cap O_E^n],\qquad J=J_\Lambda^r .
\]
This is the source index
\((\Lambda.L_n.\mathrm{addSubgroupOf}\,\Lambda.La).\mathrm{index}\).
No spanning condition on \(x_0\), or triviality of \(\Lambda\) at
the additional integer \(\mathrm{badC}(Y)\), is needed for this
norm implication; those are separate premises of \(\mathrm{thFam}\).

\begin{srproposition}\label{prop:sr-source-raw-sup-9000}
Under these exact hypotheses, for every coefficient-field
embedding \(\sigma\) use the source good embedding
\(\sigma^+=\mathrm{goodOf}(c,\sigma)\) and its actual
uniformization \(U_{\sigma^+}=Y.U(\sigma^+)\).
For every frame \(z\) in its frame domain,
\[
 \mathrm{homN}_{U_{\sigma^+},\xi,\mathrm{wtZ}(a,n)}(y;z)
 \le \sqrt J\,\rho\,
                  \exp\!\bigl(9000(frL_b)^2\kappa\bigr).
 \tag{RN3}\label{eq:sr-source-raw-sup-9000}
\]
The same constant \(9000\) works before all of the displayed
packet and family choices.
\end{srproposition}

\subsection{The elementary logarithmic estimate}
We first give the numerical producer used by the source.
Let \(f,r,n\) be positive integers with \(r\le n\), let \(d\ge0\)
be an integer, and let \(L\ge2\), \(b\ge d\), \(t,\lambda>0\),
\(\mathrm{rd}\ge1\).  Suppose
\[
 \log(1+b)\le2L,\quad \log n\le L,\quad \log f\le2L,\quad
 \log\mathrm{rd}\le2L,\quad
 \log(\lambda^{-2})\le80L.
\]
Put
\[
 X=\frac{4\sqrt{2frn}\,(2f^{3/2}n^2\mathrm{rd})}
                      {\sqrt\pi\,\lambda},\qquad
 E_d=(\pi e t d)^{d/2},
\]
using \(E_0=1\).
Let \(c_0,p_1,S,v\ge0\), \(\ell\ge0\), with
\[
 c_0\le e^{frnL},\quad p_1\le e^\ell,\quad
 S\le2v,\quad
 v\le c_0p_1E_d\sqrt2(1+X)^{frn}.
\]
Then
\[
 \log\!\left(e^{-(d/2)\log(\pi t)}S\right)
                   \le \ell+10Lb+200frnL .
 \tag{RN4}\label{eq:sr-source-log-sup}
\]
For \(S=0\), the source convention \(\log0=0\) makes the assertion
immediate, since its right side is nonnegative.
For \(S>0\), all factors on which logarithms are taken below
are positive.
The exact expansion is
\[
 \log X=\tfrac72\log2+2\log f+\tfrac12\log r+
     \tfrac52\log n+\log\mathrm{rd}
                       -\tfrac12\log\pi-\log\lambda
                \le\tfrac72\log2+49L.
\]
For \(x>0\), \(\log x\le B\), \(B\ge0\), the two cases
\(x\le1\) and \(x>1\) give
\(\log(1+x)\le\log2+B\).
Consequently \(\log(1+X)\le55L\).
If \(d>0\),
\[
 \log E_d-\tfrac d2\log(\pi t)
       =\tfrac d2(1+\log d)\le10Lb;
\]
use \(\log d\le\log(1+b)\le2L\), \(d\le b\), and \(L\ge2\).
For \(d=0\) the left side is zero.
Taking logarithms of \(S\le2v\) and the bound for \(v\),
then using
\(\log c_0\le frnL\), \(\log p_1\le\ell\), and
\(frnL\ge2\), gives (RN4); the constant terms
\(\frac32\log2\) and \(56frnL\) are bounded by \(200frnL\).
This is the direct calculation of
\srcid{LogExpMulLeOfGeneratorSupBound}, including its zero case.

\subsection{Weights, a half maximum, and the first minimum}
\begin{proof}[Proof of Proposition~\ref{prop:sr-source-raw-sup-9000}]
The actual \(\mathrm{LawsJ}\), \(\mathrm{InR}\), and
\(\mathrm{Good}\) witnesses give the source
\srcid{laws_galois_pull} on this same \(Y\).
Its relevant fields are:
\(\mathrm{IsUnf}\) for each good embedding,
\(D_w.p+D_w.q=r\), \(l_2.A=1\), the scalar
\(|\tau(\xi)|=\xi_R(D)\), and trivial arithmetic character
for these dominant even weights.
These are the precise uniformization fields used below; this
section does not reconstruct their model producer.

Pair the embeddings of \(E\) by complex conjugation.
Their CM type has \(f\) elements, and the two block
multiplicities add to \(r\).  The positive rows of
\(\mathrm{laOf}\) have absolute values \(a(\tau,i)+n/2\);
the other rows have absolute values
\(a(\bar\tau,i)+n/2\), in the source's reversed row order.
Summing these terms gives exactly (RN1).
It follows that
\[
 d\le\kappa,\quad frn\le2\kappa,\quad
 n\le2\kappa,\quad \kappa>0.
 \tag{RN5}\label{eq:sr-source-raw-size-inequalities}
\]
For each embedding the two endpoints bounding all rows differ
by at most
\(\sum_i a(\tau,i)+\sum_i a(\bar\tau,i)+n\).
The source CM-type relabelling is bijective.  Summing that bound
therefore gives the actual spread estimate
\[
 \sum_\tau\mathrm{spread}(\lambda;\tau)\le d+frn\le3\kappa .
 \tag{RN6}\label{eq:sr-source-raw-spread}
\]
This is \srcid{SumWeightSpreadLeExponentNormAdd.sum_spread_laU_le}.

Let \(S\) be the finite supremum of the actual weighted frame
norm of the reading of \(y\) in \(U_{\sigma^+}\).
Its finiteness is the \(\mathrm{holSub}\) conclusion in the
definition of this cusp space.
Every weighted point norm is at most \(S\).
If \(S=0\), (RN3) follows immediately.
Assume \(S>0\).
Apply Proposition~\ref{prop:sr-source-orthogonal-halfmax},
with \(w=3\kappa\), to get an orthogonal frame \(x_f\) in one
actual component, with
\[
 S\le2|\mathrm{ev}_{\sigma^+}(y;x_f)|,\qquad
 \mathrm{firstMinimumSq}(x_f)
       \ge2\epsilon_{\rm lev}(3\kappa)/\mathrm{rd}_E^2.
 \tag{RN7}\label{eq:sr-source-raw-halfmax}
\]
This uses a half supremum and the actual Borel
orthogonalization; it does not require a maximum.

Put \(u=e^{L_b}\ge2\).  The packet bounds give
\[
 M\le u^{30},\quad \mathrm{rd}_E\le u^3,\quad
 \sqrt{D.\mathrm{host}.d}\le u^3,\quad
 f\le u,\quad 3\kappa\le3u^{20}.
\]
For the middle host bound use the source arithmetic comparison
\(\log(D.\mathrm{host}.d)\le2\log\mathrm{rd}_E\).
Here \(\log\mathrm{rd}_E\le\mathrm{fieldSz}(E)\le3L_b\).
Thus \(f(3\kappa)+2\le u^{24}\), and the exact denominator in
\(\epsilon_{\rm lev}\) satisfies
\[
 2M\sqrt{D.\mathrm{host}.d}\,(3f\kappa+2)\log2
       \le2u^{57}\le\pi u^{57}.
\]
Both entries in the minimum defining \(\epsilon_{\rm lev}\)
are at least \(u^{-57}\).
Consequently
\[
 \epsilon_{\rm lev}(3\kappa)\ge e^{-57L_b},\qquad
 \lambda=e^{-40L_b}\ \Longrightarrow\
      \lambda^2\le\mathrm{firstMinimumSq}(x_f).
 \tag{RN8}\label{eq:sr-source-raw-lambda}
\]
Indeed \(\lambda^2\mathrm{rd}_E^2\le e^{-74L_b}
\le e^{-57L_b}\le2\epsilon_{\rm lev}(3\kappa)\).
This is the source's \srcid{inv_pow_le_epStar} substitution.

\subsection{The finite coefficient sum and the absolute constant}
The harmonic polynomial has degree \(d\).
There are \(frn\) letters: pairing conjugate embeddings makes
the total block multiplicity \(fr\), with \(n\) columns.
Every support exponent is between \(0\) and \(d\).
Writing \(N_P=|\operatorname{supp}P|\) gives
\[
 N_P\le(d+1)^{frn}.
\]
For the fixed good embedding put
\(p_P=\sum_\alpha|\sigma^+(P_\alpha)|\).
Cauchy--Schwarz and the definition of the polynomial norm give
\[
 p_P\le\sqrt{N_P}\,\mathrm{polyNorm}_c(P)
                                      \le\sqrt{N_P}\,\rho .
 \tag{RN9}\label{eq:sr-source-raw-polynomial-sum}
\]
The polynomial norm is the square root of the sum of coefficient
squares over all coefficient-field embeddings, so it bounds the
sum of squares at this one embedding.
In particular \(\rho\ge0\).

The phase of \(\phiHat(\mathrm{indPhi}(x_0))\) has modulus one
at every coefficient-field embedding.
Indeed \(\mathrm{indPhi}\) has exactly one nonzero term,
and its Fourier value is an integral power of the specified
root of unity.  This root has nonzero order
\(l_2.M'=\mathrm{lcm}(1,l_2.M)\), because \(l_2.A=1\).
Its image under any embedding has modulus one, by raising
to that nonzero order.
This checks the coefficient-function premise of the raw-value
bound without any size assumption on \(x_0\).

Use (RN8) in the exact raw-value estimate
\eqref{eq:sr-source-raw-value-bound}, with \(t=2\xi\).
For \(v=|\mathrm{ev}_{\sigma^+}(y;x_f)|\) it reads
\[
 v\le\sqrt J\,p_P(\pi e t d)^{d/2}\sqrt2
 \left(1+\frac{4\sqrt{2frn}\,(2f^{3/2}n^2\mathrm{rd}_E)}
                        {\sqrt\pi\,e^{-40L_b}}\right)^{frn}.
 \tag{RN10}\label{eq:sr-source-raw-numeric-input}
\]
The same actual connection, frame and all-test-ring equality
(RN2) supply this formula.

We now apply (RN4) with
\[
 L'=21L_b,\quad b=\kappa,\quad \ell=0,\quad
 c_0=p_1=1,\quad
 S'=S/(\sqrt J\,p_P),\quad v'=v/(\sqrt J\,p_P).
\]
The division is legitimate:
\(S>0\), \(S\le2v\), and (RN10) force \(\sqrt J\,p_P>0\).
By (RN5), \(n\le2\kappa\); hence
\[
 \log n\le21L_b,\qquad
 \log(1+\kappa)\le22L_b,\qquad
 \log f\le L_b,\qquad
 \log\mathrm{rd}_E\le3L_b.
\]
For \(\log(1+\kappa)\), use \(\kappa\ge1/2\) from
\(frn\ge1\), so \(1+\kappa\le3\kappa\).
These bounds verify every logarithmic premise at \(L'\).
Also \(\log(\lambda^{-2})=80L_b\le80L'\).
The conclusion of (RN4), exponentiated and rescaled, gives
\[
 e^{-(d/2)\log(\pi t)}S
 \le\sqrt J\,p_P
       \exp\!\left(210L_b\kappa+4200frnL_b\right).
\]
Next (RN5) and \(\log(d+1)\le22L_b\) imply
\(\sqrt{N_P}\le e^{22L_b\kappa}\).
Using (RN9), then \(frn\le2\kappa\), bounds the preceding
right side by
\[
 \sqrt J\,\rho\,
 \exp\!\left((210+8400+22)L_b\kappa\right)
 \le\sqrt J\,\rho\,
                    e^{9000(frL_b)^2\kappa}.
 \tag{RN11}\label{eq:sr-source-raw-final-number}
\]
The last step uses \(L_b\le(frL_b)^2\).

The exact source Lie factor is
\[
 \mathrm{lieFacZ}(D_w,\xi,\mathrm{wtZ}(a,n))
                 =e^{-(d/2)\log(\pi(2\xi))}.
\]
Multiplying the point-norm bound by this factor, and using
\(t=2\xi\), turns (RN11) into (RN3).
This proves the original constant \(9000\).
\end{proof}

\begin{srcorollary}\label{cor:sr-source-smul-raw-sup-9000}
Retain all preceding data.
Let \(s\in c.F\), \(B\ge0\), and suppose that for every
coefficient-field embedding \(\sigma\),
\[
                         |\sigma(s)|^2J\le B^2 .
 \tag{RN12}\label{eq:sr-source-raw-scalar-bound}
\]
Then, in every indicated domain frame,
\[
 \mathrm{homN}_{U_{\sigma^+},\xi,\mathrm{wtZ}(a,n)}(s\,y;z)
              \le B\rho\,e^{9000(frL_b)^2\kappa}.
 \tag{RN13}\label{eq:sr-source-smul-raw-sup-9000}
\]
\end{srcorollary}
\begin{proof}
The actual reading is semilinear and its frame norm is
absolutely homogeneous.  Thus
\(\mathrm{homN}(s\,y;z)=|\sigma^+(s)|\mathrm{homN}(y;z)\).
Both sides of
\(|\sigma^+(s)|\sqrt J\le B\) are nonnegative, and their
squares are ordered by (RN12).
Multiplying (RN3) by \(|\sigma^+(s)|\) proves (RN13).
This is \srcid{homN_smul}, followed by
\srcid{homN_smul_raw_le}; its scalar bound is index-weighted.
\end{proof}

\paragraph{Frozen-source correspondence and scope.}
The numerical estimate (RN4) is
\srcid{src/Development/AbelianVarietyTheory/Proofs/CuspFormEvaluations/LogExpMulLeOfGeneratorSupBound.lean},
blob \srcid{89ce9b945e9447d0ed1d20c550115b026e2dd8bb}.
The spread calculation is
\srcid{src/Development/AbelianVarietyTheory/Proofs/DivisorProofs/SumWeightSpreadLeExponentNormAdd.lean}.
The assembly, including \srcid{sizeK_eq}, \srcid{card_letters},
\srcid{sum_coeff_le}, \srcid{inv_pow_le_epStar},
\srcid{numbers}, \srcid{norm_phiHat_indPhi}, and both final
norm theorems, is
\srcid{src/Development/AbelianVarietyTheory/Proofs/DivisorProofs/NormPhiHatIndPhi.lean},
blob \srcid{b7fb0bf05a120ad2a79381a586e429c8525f82b2}.
All are read at \srcid{alpoge-source-7cbf7a0c}.

The antecedents retain the actual deep-cover Galois laws,
the same-model reading and connection, and the geometric
analytic-fibre/level/vacuum comparisons used by the preceding
raw-value theorem.  This section proves their exact norm
implication and its uniform numerical constant.
It does not assert that these geometric suppliers have been
constructed by the numerical estimates, or that this norm
bound alone proves spanning.
